% !PRINTED 124
\weberpage{124}
%<<CONT
and one then sees that these elements are all different from the elements $C$, since, if $g'^{-1}c'g' = g^{-1}cg$ were the case, we should also have $c' = g'g^{-1}cgg'^{-1}$, and $c'$ would thus, contrary to the assumption, already occur in $C$. In this way one continues with the formation of the groups $B, B', B'', \ldots$ and the systems $C, C', C'', \ldots$ until all the elements of $P$ are accommodated in the systems $A, C, C', C'', \ldots$. One then has the numerical equations:
\weq{9}{n = \nu j = \mu\varepsilon = \mu'\varepsilon' = \mu''\varepsilon''\ldots}
\weq{10}{n = \nu + \varepsilon + \varepsilon' + \varepsilon'' + \cdots}

On the basis of these formulae \gesperrt{Sylow}'s theorem can now be proved by complete induction, in which two cases are to be distinguished:

1) \gesperrt{The degree} $\nu$ \gesperrt{of} $A$ \gesperrt{is divisible by} $p$\gesperrt{.} In this case, since $A$ is an Abelian group, there is by \S.\,9, 2. an element $a$ in $A$ of degree $p$, and the cyclic group of degree $p$
\[ 1,\ a,\ a^2,\ \ldots,\ a^{p-1}, \]
which we denote by $Q$, is a normal divisor of $P$, because for an element $a$ of $A$ we should always have $x^{-1}ax = a$. The complementary group to $Q$, $P/Q$ is of degree $n:p$, and its degree is therefore smaller than $n$ and divisible by $p^{\varkappa-1}$. By our assumption there is therefore a divisor of $P/Q$ of degree $p^{\varkappa-1}$, and consequently by \S.\,6, 2. there is a divisor of $P$ of degree $p^{\varkappa}$.

2) \gesperrt{The degree} $\nu$ \gesperrt{of} $A$ \gesperrt{is not divisible by} $p$\gesperrt{.} Since $n$ is divisible by $p$, in this case by (10) the numbers $\varepsilon, \varepsilon', \varepsilon'', \ldots$, which are all greater than 1, cannot all be divisible by $p$. Now if $\varepsilon$ is not divisible by $p$, then $\mu$ is divisible by $p^{\varkappa}$. The degree of the group $B$ is smaller than $n$ and divisible by $p^{\varkappa}$, and so by assumption there is a divisor of $B$, which is therefore also a divisor of $P$, of degree $p^{\varkappa}$.

With this Theorem I is proved.

To this theorem, however, there are still essential supplements to be added.
%CONT>>
